\section{Related work}
\label{sec:related}

\paragraph{Music representation models}
MERT~\cite{li2024mert} is a BERT-style transformer trained with a multi-task
masked-prediction objective guided by a residual vector-quantisation acoustic
teacher and a constant-Q transform musical teacher. It reaches competitive
results across fourteen music information retrieval tasks while remaining
computationally modest, and is released as a reusable feature extractor rather
than a detector, which makes it a natural frozen backbone. Layer-wise analyses
of self-supervised audio models show that depth encodes a transition from
low-level acoustic to higher-level semantic features~\cite{elkheir2025}, which
motivates learned aggregation across layers rather than use of a single layer.

\paragraph{AI-music and synthetic-song detection}
Cros Vila et al.~\cite{crosvila2025} establish the cross-generator protocol and
the $0.99 \rightarrow 0.629$ collapse this work takes as its reference point.
SONICS~\cite{rahman2025sonics} supplies over 97{,}000 songs of real and
end-to-end generated music from Suno and Udio, together with a spectro-temporal
tokeniser that improves long-range modelling. Afchar et al.~\cite{afchar2025icassp}
reframe detection as identifying a neural decoder's fingerprint and report 99.8\%
accuracy while showing candidly that the accuracy is fragile under distribution
shift. FakeMusicCaps~\cite{comanducci2025} extends the problem to attribution
across five text-to-music generators and reports strong closed-set but weak
open-set performance. A frequency-domain analysis~\cite{afchar2025ismir} grounds
detectability in interpretable spectral evidence.

\paragraph{Shortcuts and evaluation validity}
The closest methodological precedent is in singing-voice deepfake detection,
where CtrSVDD~\cite{zang2024ctrsvdd} was constructed specifically to isolate
synthesis artefacts from voice-separation artefacts, and SingFake~\cite{zang2024singfake}
showed that speech-trained models degrade sharply under unseen singers and
codecs. Both warn that pre-processing can confound a classifier. Our shortcut
and bandwidth diagnostics extend that warning to the music setting and make it
measurable.

\paragraph{Explainability for audio detection}
Attention rollout for transformer audio classifiers~\cite{channing2025} frames
transparency as a prerequisite for deployment. Time-domain relevance
attribution~\cite{grinberg2025} identifies which temporal segments drive a
decision, benchmarked against Grad-CAM and SHAP baselines. SHAP has been applied
to spectrogram- and waveform-trained spoofing networks~\cite{yu2024}, and
Mel-spectrogram CNNs with Grad-CAM-style localisation have been used to
highlight fake-indicative regions~\cite{bajwa2025}. These methods target speech
almost exclusively; an early benchmark for explainable machine-generated music
detection~\cite{li2025explainable} notes that explanation fidelity remains an
open metric. Our results bear directly on that: we find temporal attribution
faithful for real audio and unfaithful for generated audio.
